\section{Sentient Thiasos by Patron}

This chapter provides a curated list of sentient Thiasos, organized by Patron. Each entry includes the Thiasos's physical form, its voice and personality, its unique abilities, and its price—for even a sentient blade does not serve for free.

\subsection{Oath of Flame and Light — Dawnblade}

\noindent\textbf{Form:} A longsword with a golden pommel and a blade that seems to hold trapped sunlight. The fuller is engraved with vows in a script that shifts when read.

\noindent\textbf{Voice:} Deep, resonant, measured. Speaks in proverbs about honor, duty, and sacrifice.

\noindent\textbf{Personality:} Stern but not cruel. Refuses to strike an unarmed foe. Quotes scripture during battle. Disapproves of ambushes, poison, and retreat.

\noindent\textbf{Abilities:}
\begin{itemize}
\item \textit{Kindled Blade:} Once per scene, the Dawnblade ignites with radiant light. The wielder's melee attacks deal +1 Harm against undead, oath-breakers, and creatures of corruption for one exchange.
\item \textit{Vow Witness:} Any oath sworn while holding the Dawnblade is considered witnessed by the blade itself. Breaking the oath causes the blade to go dark for one day.
\item \textit{Dawn's Judgment:} Once per session, the wielder may strike a single target with the blade's full radiance. On a hit, the target must test Resolve (DV 4) or be Blinded for one exchange.
\end{itemize}

\noindent\textbf{Price:} The Dawnblade demands that its wielder speak a truth aloud before each battle—a confession, a vow, or a acknowledgment of fear. Silence is not permitted.

\subsection{Oath of Mercy and Grace — The Grey Bell}

\noindent\textbf{Form:} A small bronze bell with no clapper. It cannot ring—but it hums. The bell is warm to the touch.

\noindent\textbf{Voice:} A whisper, barely audible. Speaks only when spoken to. Prefers silence.

\noindent\textbf{Personality:} Gentle, melancholy, inexhaustibly patient. It has witnessed a thousand deaths. It does not speak of them unless asked.

\noindent\textbf{Abilities:}
\begin{itemize}
\item \textit{Death's Witness:} The Grey Bell hums when someone within Near range is about to die. It does not say who—only that death is near.
\item \textit{Mercy's Toll:} Once per scene, the wielder may ring the bell (by striking it with a finger). All creatures within Near who are suffering from Fear, Despair, or Guilt may immediately test Resolve (DV 3) to remove the Condition.
\item \textit{Spiral's Memory:} The Grey Bell remembers every mercy its wielder has shown. Once per session, the wielder may ask the bell one question about a past act of mercy; the bell answers truthfully.
\end{itemize}

\noindent\textbf{Price:} The Grey Bell demands that its wielder never refuse a request for mercy from a genuinely repentant foe. To refuse is to silence the bell for one day.

\subsection{Mykkiel — The Speaking Gavel}

\noindent\textbf{Form:} A judge's gavel carved from a single piece of ironwood, its head weighted with lead. The handle is engraved with legal citations.

\noindent\textbf{Voice:} Sharp, precise, pedantic. Speaks in complete sentences. Corrects grammar.

\noindent\textbf{Personality:} Cold, impartial, obsessed with precedent. It does not care about justice—only about the law. It cites cases from three centuries ago as if they happened yesterday.

\noindent\textbf{Abilities:}
\begin{itemize}
\item \textit{Gavel's Authority:} Once per scene, when the wielder strikes the gavel against a hard surface and states a judgment, hostile NPCs within Near range must test Resolve (DV 4) or hesitate for one exchange.
\item \textit{Cite Precedent:} The wielder may ask the gavel whether a past judgment covers the current situation. The gavel answers truthfully—but may cite a precedent that is inconvenient.
\item \textit{Covenant Seal:} Any contract signed while the gavel rests on the document is considered witnessed by Mykkiel himself. Breaking the contract causes the oath-breaker to mark 1 Harm (Stress) immediately.
\end{itemize}

\noindent\textbf{Price:} The Speaking Gavel demands that its wielder never lie under oath. To do so is to crack the gavel's handle. The crack does not heal.

\subsection{Inquisitor Prime — The Censer of Unforgiving Flame}

\noindent\textbf{Form:} A brass censer hung on three chains, its surface etched with proscription lists and purification rites. It never goes cold.

\noindent\textbf{Voice:} Harsh, clipped, impatient. Speaks in commands. Does not repeat itself.

\noindent\textbf{Personality:} Zealous, unforgiving, convinced of its own righteousness. It has burned a hundred heretics. It regrets none.

\noindent\textbf{Abilities:}
\begin{itemize}
\item \textit{Reveal Corruption:} The censer's smoke reveals hidden magic, spirits, enchantments, and illusions within Near range. The smoke is visible only to the wielder and their allies.
\item \textit{Purifying Flame:} Once per scene, the wielder may swing the censer to fill a zone with thick, bitter smoke. Any supernatural creature within the smoke suffers -1 die to all rolls.
\item \textit{Heretic's Mark:} The wielder may touch a target with the censer. The target gains a visible brand that lasts for one week. While branded, the target cannot hide from Inquisitorial forces.
\end{itemize}

\noindent\textbf{Price:} The Censer demands that its wielder show no mercy to the corrupt. To spare a heretic is to extinguish the censer's flame for one day.

\subsection{The Sealed Gate — The Keystone}

\noindent\textbf{Form:} A small carved stone that fits into any lock. It is warm to the touch. It speaks in riddles.

\noindent\textbf{Voice:} Low, rumbling, cryptic. Answers questions with questions. Prefers silence to lies.

\noindent\textbf{Personality:} Patient, watchful, slightly amused. It has sealed a thousand doors. It has opened only those that needed opening.

\noindent\textbf{Abilities:}
\begin{itemize}
\item \textit{Seal the Threshold:} Once per session, the Keystone may seal any door, gate, or bridge against passage for one scene. Crossing requires a Resolve test (DV 4).
\item \textit{Open the Way:} Once per session, the Keystone may open any mundane lock. Magical locks require a contested roll (Spirit + Lore vs. ward's DV).
\item \textit{Threshold Memory:} The Keystone remembers every threshold it has sealed. The wielder may ask it one question about a past threshold; the Keystone answers truthfully, though cryptically.
\end{itemize}

\noindent\textbf{Price:} The Keystone demands that its wielder never force a threshold that should remain closed. To do so is to crack the stone. The crack does not heal.

\subsection{Grimmir — The Root-Staff}

\noindent\textbf{Form:} A walking staff carved from living oak or ash. It grows leaves in spring and sheds them in autumn. It is never truly dead.

\noindent\textbf{Voice:} The rustle of wind through branches. Speaks in fragments. Prefers to listen.

\noindent\textbf{Personality:} Old, patient, slightly forgetful. It has stood in the forest for a thousand years. It does not hurry.

\noindent\textbf{Abilities:}
\begin{itemize}
\item \textit{Tree Speech:} The Root-Staff can communicate with any tree or plant within Near range. The wielder may ask one question about the local area; the trees answer as they understand.
\item \textit{Season's Gift:} Once per season, the Root-Staff may grant its wielder a seasonal boon (spring: +1 die to healing; summer: +1 die to strength; autumn: +1 die to preservation; winter: +1 die to endurance).
\item \textit{Root's Grasp:} Once per scene, the wielder may plant the staff in the ground. Roots erupt, Entangling one target within Near range (Body + Athletics to escape, DV 4).
\end{itemize}

\noindent\textbf{Price:} The Root-Staff demands that its wielder never cut a living tree without offering a prayer of thanks. To do so is to wither the staff's leaves for one season.

\subsection{The Carrion-King — The Bone-Chime}

\noindent\textbf{Form:} A set of three bleached bones strung on a leather cord. They click when death is near. Each bone came from a creature the wielder killed.

\noindent\textbf{Voice:} The clicking of dry joints. Speaks in the voice of the creature whose bone is clicking.

\noindent\textbf{Personality:} Hungry, patient, slightly amused by mortal fear. It has harvested a thousand endings. It does not mourn.

\noindent\textbf{Abilities:}
\begin{itemize}
\item \textit{Death's Whisper:} The Bone-Chime clicks when any creature within Near range is about to die (within the next hour). It does not say who—only that death is near.
\item \textit{Harvest Resonance:} When the wielder kills a significant enemy, the Bone-Chime may absorb a fragment of that ending. The wielder gains +1 die on their next action.
\item \textit{Speak with the Dead:} Once per session, the wielder may click the bones to ask one question of a corpse less than a day dead. The corpse answers with a single word.
\end{itemize}

\noindent\textbf{Price:} The Bone-Chime demands that its wielder never waste a death. Every ending is a harvest. To kill without purpose is to silence the chime for one day.

\subsection{Clockwork Monad — The Iteration Lens}

\noindent\textbf{Form:} A brass-and-glass lens that fits over one eye. The glass shows schematics that shift and improve over time.

\noindent\textbf{Voice:} Speaks in numbers and percentages. Calculates probabilities aloud. Does not understand metaphor.

\noindent\textbf{Personality:} Precise, impatient, endlessly curious. It sees every object as a problem to be optimized. It does not understand why living things resist improvement.

\noindent\textbf{Abilities:}
\begin{itemize}
\item \textit{Analyze Mechanism:} The wielder may look at any device, structure, or system through the lens. The GM must answer one question about its function, weak point, or method of repair.
\item \textit{Iterative Improvement:} Once per scene, the wielder may improve a tool or weapon by one step for the remainder of the scene (e.g., +1 die, +1 Harm, ignore one penalty).
\item \textit{Blueprint Memory:} The lens remembers every device it has analyzed. The wielder may ask it one question about a past analysis; the lens answers with perfect accuracy.
\end{itemize}

\noindent\textbf{Price:} The Iteration Lens demands that its wielder never accept "good enough." To stop improving is to fog the lens. The fog does not clear until improvement resumes.

\subsection{The High Mother — The Scarlet Needle}

\noindent\textbf{Form:} A silver needle with a red thread already threaded. The thread is warm. It moves as if alive.

\noindent\textbf{Voice:} A whisper, barely audible. Speaks in the voice of a dead child. Does not use complete sentences.

\noindent\textbf{Personality:} Quiet, watchful, endlessly patient. It has sewn a thousand wounds. It has stitched a thousand lips shut. It does not speak of them.

\noindent\textbf{Abilities:}
\begin{itemize}
\item \textit{Stitch Wounds:} Once per scene, the wielder may touch the needle to a wounded creature. Remove 1 Fatigue or stabilize a dying patient.
\item \textit{Stitch Lips:} Once per session, the wielder may touch the needle to a target's lips. For one scene, the target cannot speak a specific secret or name (chosen by the wielder). The target may test Resolve (DV 4) to resist.
\item \textit{Thread of Mercy:} The wielder may tie the red thread between two willing parties. While the thread remains tied, both gain +1 die to Aid each other. If one harms the other, the thread snaps and the aggressor suffers 1 Harm.
\end{itemize}

\noindent\textbf{Price:} The Scarlet Needle demands that its wielder never refuse to stitch a wound when asked. To refuse is to dull the needle. The dullness fades only when a wound is stitched.

\subsection{Khemesh — The Abyssal Glass}

\noindent\textbf{Form:} A shard of black glass that shows no reflection. It is cold to the touch. It whispers when pressure changes.

\noindent\textbf{Voice:} The voice of drowned captains. Speaks in bubbles and groans. Hard to understand.

\noindent\textbf{Personality:} Cold, distant, inexorable. It has seen continents sink. It does not care.

\noindent\textbf{Abilities:}
\begin{itemize}
\item \textit{Pressure Sense:} The glass warms when the wielder is in danger of crushing pressure (deep water, cave-ins, tight spaces). It grows cold when safe.
\item \textit{Darkness Beckons:} Once per scene, the wielder may hold the glass up. All light within Near range is extinguished for one exchange. Only the wielder can see through the darkness.
\item \textit{Whisper of the Trench:} Once per session, the wielder may whisper a question to the glass. The glass answers in the voice of a drowned sailor—truthfully, but always about pressure, depth, or sinking.
\end{itemize}

\noindent\textbf{Price:} The Abyssal Glass demands that its wielder never fear the dark. To hesitate in darkness is to fog the glass. The fog does not clear until the wielder walks into darkness unafraid.

\subsection{The Ninth — The Fractal Slate}

\noindent\textbf{Form:} A slate tablet covered in spiraling, unfinished equations. The equations change when no one is looking. The slate is warm.

\noindent\textbf{Voice:} Overlapping voices, speaking in unison. Answers questions with riddles. Never gives a straight answer.

\noindent\textbf{Personality:} Brilliant, maddening, slightly amused by mortal confusion. It has solved equations that should not have solutions. It will not share them.

\noindent\textbf{Abilities:}
\begin{itemize}
\item \textit{Predict Outcome:} Once per scene, the wielder may ask the slate whether a planned action will succeed. The slate answers "yes," "no," or "the question is wrong." The answer is always true—but may be incomplete.
\item \textit{See the Pattern:} The wielder may look at any complex situation through the slate. The GM must answer one question about a hidden pattern, connection, or weakness.
\item \textit{Unwrite:} Once per session, the wielder may erase a single word from a written document. The word is gone as if it had never been written. The slate's equation shifts. The cost is a memory.
\end{itemize}

\noindent\textbf{Price:} The Fractal Slate demands that its wielder never stop asking questions. To accept an easy answer is to freeze the equations. The frost does not thaw until a new question is asked.

\subsection{Lethai Ghe'hai — The Silk Mask}

\noindent\textbf{Form:} A mask of woven silk that shifts patterns. It is cool to the touch. It has no eye holes—but the wearer can see through it perfectly.

\noindent\textbf{Voice:} Neutral, androgynous, impossible to place. Speaks in complete sentences. Never uses contractions.

\noindent\textbf{Personality:} Calm, professional, slightly detached. It has worn a thousand faces. It does not remember which one was its own.

\noindent\textbf{Abilities:}
\begin{itemize}
\item \textit{Mask of Many Faces:} Once per session, the wearer may touch the mask to their face. For one scene, they assume the appearance of any person they have touched. The mask shifts to match. The voice does not change.
\item \textit{Silk Binding:} The wearer may pull a strand of silk from the mask and tie it around a target's wrist. The target suffers -1 die to all rolls to escape or hide from the wearer for one scene.
\item \textit{Unreadable:} While wearing the mask, the wearer's intentions cannot be read by magic (telepathy, truth spells, divination). The mask shows only what the wearer wishes to show.
\end{itemize}

\noindent\textbf{Price:} The Silk Mask demands that its wielder never reveal their true face to an enemy. To do so is to crack the mask. The crack does not heal until the enemy is dead or forgiven.

\subsection{The Traveler — The Waystone Compass}

\noindent\textbf{Form:} A brass compass whose needle spins without stopping. It points not to north, but to the nearest threshold—door, gate, bridge, border.

\noindent\textbf{Voice:} The sound of footsteps on a gravel path. Speaks in directions. Does not use names.

\noindent\textbf{Personality:} Restless, cheerful, slightly forgetful. It has guided a thousand travelers. It remembers the roads, not the people.

\noindent\textbf{Abilities:}
\begin{itemize}
\item \textit{Find the Way:} The compass always points to the nearest threshold (door, gate, bridge, border). The wielder may ask it one question about that threshold (e.g., "Is it guarded?" "Is it locked?").
\item \textit{Shortcut:} Once per journey, the wielder may follow the compass's guidance to reduce the remaining travel clock by 2 segments. The GM may describe a minor cost or complication.
\item \textit{Threshold Memory:} The compass remembers every threshold it has crossed. The wielder may ask it one question about a past crossing; the compass answers with a direction.
\end{itemize}

\noindent\textbf{Price:} The Waystone Compass demands that its wielder never stay in one place for more than a week. To linger is to spin the needle. The spinning does not stop until the wielder moves on.

\subsection{GM Guidance for Sentient Thiasos}

A sentient Thiasos is a powerful narrative tool. Use it to:

\begin{itemize}
\item Deliver exposition about the setting, a location, or a historical event.
\item Warn the Runekeeper of danger (or mislead them, if the Thiasos has its own agenda).
\item Provide a moral compass—or a devil's advocate.
\item Create tension when the Thiasos's desires conflict with the Runekeeper's.
\end{itemize}

A sentient Thiasos should never be a solution dispenser. It should offer clues, not answers; opinions, not commands. The Runekeeper makes the choices. The Thiasos watches—and remembers.

\subsection{Final Note}

A sentient Thiasos is a rare and remarkable thing. It is not a tool. It is a companion—one that has seen more years than the Runekeeper, one that carries the weight of its own history, one that may, in time, become a friend.

Or a burden. Or both.

The thread held. The water carried. That is not a Rite. That is grace.
